% -------------------------------------------------
% Anexo E. Elementos de gamificación implementados en el prototipo
% -------------------------------------------------
\section*{Anexo E. Elementos de gamificación implementados en el prototipo}
\addcontentsline{toc}{section}{Anexo E. Elementos de gamificación implementados en el prototipo}
\label{anexo:gamificacion}

La siguiente tabla presenta una comparación entre los elementos de gamificación descritos por Watini y Setyowati (2023) y su implementación dentro de los módulos \textit{Aprender a Aprender} y \textit{Escucha Activa} del prototipo desarrollado. Esta tabla complementa la descripción teórica incluida en el capítulo correspondiente al objetivo específico 3 y permite visualizar, de manera sintética, la relación entre la fundamentación académica y las funcionalidades implementadas.

\vspace{0.4cm}

\begin{table}[H]
\centering
\footnotesize
\setlength{\tabcolsep}{6pt}
\renewcommand{\arraystretch}{1.4}
\begin{tabular}{|p{3.2cm}|p{3.2cm}|p{4.2cm}|p{4.2cm}|}
\hline
\textbf{Elemento de gamificación} & \textbf{Fundamento (Watini \& Setyowati, 2023)} & \textbf{Implementación en Aprender a Aprender} & \textbf{Implementación en Escucha Activa} \\
\hline

\textbf{Progresión por niveles} &
Progresión incremental mediante unidades pequeñas de avance. &
Estructura metas–objetivos–tareas que permite avanzar gradualmente. &
Rondas secuenciales que se desbloquean al completar la anterior. \\
\hline

\textbf{Retroalimentación inmediata} &
La motivación aumenta con feedback instantáneo tras una acción. &
Autoevaluaciones por tarea; actualización del estado en el Kanban. &
Puntuación inmediata por ronda; comparación entre respuestas y criterios. \\
\hline

\textbf{Indicadores visuales de progreso} &
El progreso visible guía el comportamiento y mantiene el compromiso. &
Estados Kanban: Por hacer, En progreso, En pausa, Completado. &
Barra de progreso por ronda y resumen final de desempeño. \\
\hline

\textbf{Dificultad ajustada} &
Los desafíos deben ser equilibrados para evitar frustración. &
Complejidad variable de objetivos y tareas según el plan del estudiante. &
Configuración del nivel de dificultad (longitud del audio, complejidad semántica, modo de juego). \\
\hline

\textbf{Personalización y agencia} &
Permitir decisiones incrementa el sentido de control y la motivación. &
Creación, edición y copia de planes; personalización del roadmap. &
Configuración de rondas, prompts, dificultad y reutilización de desafíos. \\
\hline

\textbf{Anticipación del siguiente paso} &
Mantener al usuario en flujo reduce el abandono. &
Indicaciones para la siguiente tarea y objetivo dentro del roadmap. &
Transiciones automáticas entre rondas y vista previa del avance. \\
\hline

\textbf{Socialización indirecta} &
La percepción de comunidad actúa como motivador. &
Exploración de planes públicos creados por otros estudiantes. &
No aplica de forma directa; el enfoque es individual. \\
\hline

\end{tabular}
\caption{Relación entre elementos de gamificación y su implementación en el prototipo.}
\label{table:gamification-elements}
\end{table}